\ncdlesson{10}{Præstationssamfundet og jagten på det perfekte liv}{绩效社会与对完美生活的追逐}
{DU3 Modul 5 (PD3) $\rightarrow$ Modul 6 (Studieprøven, B2 $\rightarrow$ C1)｜对标 NCE 4 \& 现代社会心理学与文化哲学评论}
{\textbf{哲学与心理学思辨评论}（连续设问推进 \textit{Hvornår blev...? Er det i grunden...?} · 概念拆解句型 \textit{Begrebet dækker ikke blot over..., men...} · 悖论揭示 \textit{Paradokset består i, at...}）· 韩炳哲（Byung-Chul Han）《倦怠社会》语汇体系}
{欢迎来到第 10 课！在高福利的丹麦，为什么越来越多年轻人仍感到喘不过气，陷入“12 分女孩（\textit{12-talspiger}）”、完美履历与职业倦怠（\textit{udbrændthed}）的焦虑漩涡？在 PD3 与 Studieprøven 中，哲学与心理反思类评论最忌空洞说教，而最推崇\textbf{以层层递进的连续设问刺破表象}，再用\textbf{严密的概念拆解与悖论揭示}展开深度推演。本课借助韩炳哲（Byung-Chul Han）的\textbf{“绩效社会”（\textit{præstationssamfund}）}理论，全面升级你的 C1 思辨笔法！}

% ---------------------------------------------------------------------
\begin{foerduleaeser}
带着下面两个问题阅读 Tekst A 和 Tekst B：
\begin{enumerate}
  \item \textit{Hvordan adskiller filosoffen Byung-Chul Han ifølge Tekst A det moderne præstationssamfund fra det traditionelle disciplinsamfund, og hvorfor skaber selvoptimering en følelse af utilstrækkelighed?}（根据 Tekst A，哲学家韩炳哲如何区分现代“绩效社会”与传统“规训社会”？为什么“自我优化”反而催生了持续的匮乏感与不足感？）
  \item \textit{Hvad mener psykologen Kasper i Tekst B med ”evalueringsfrie rum”, og hvorfor påpeger Maja, at individet ikke kan løse præstationspresset alene?}（在 Tekst B 中，心理学家卡斯珀所说的“免于评价的空间”指什么？为什么玛雅指出仅靠个人无法独自化解绩效压力？）
\end{enumerate}
\end{foerduleaeser}

% ---------------------------------------------------------------------
\ncdsubsec{二、Tekster（哲学思辨专栏与校园心理对话双语篇）}

\subsection*{Tekst A：\textit{Når frihed bliver et nyt fængsel}（哲学与心理学思辨长文 · 约 370 词）}

\begin{tipbox}{语篇说明}
本篇是一篇极具思辨锋芒的文化哲学评论。阅读时请重点标出\textbf{\color{NCDNavy}深蓝色标注的连续设问推进与概念界定框架}，以及\textbf{\color{NCDRed}深红色标注的绩效社会核心概念与悖论转折词}。
\end{tipbox}

\da{Hvornår blev} hverdagen et evigt projekt, der hele tiden skal forbedres? \da{Hvornår blev} fritiden reduceret til endnu en arena for optimering af krop, netværk og CV? \da{Er det i grunden} ægte frihed, vi oplever i dag, eller blot en mere subtil form for indre tvang? I et velfærdssamfund som det danske, hvor de fleste materielle grundbehov er dækket, burde roen egentlig sænke sig. Alligevel vokser antallet af unge og voksne, der rammes af stress, søvnløshed og dyb \dabold{udbrændthed}.

For at forstå denne gåde peger den sydkoreansk-tyske filosof Byung-Chul Han på et afgørende historisk skift. I det klassiske \dabold{disciplinsamfund} blev mennesket styret udefra gennem forbud, autoriteter og klare pligter: ”Du må ikke” og ”Du skal”. Det moderne \dabold{præstationssamfund} taler derimod med en tilsyneladende frigørende stemme: ”Du kan alt!” \da{Begrebet dækker ikke blot over} stigende krav på arbejdsmarkedet eller jagten på topkarakterer i gymnasiet, \da{men over} en gennemgribende mentalitet, hvor individet gør sig selv til både entreprenør og arbejder i sit eget liv.

\da{Paradokset består i, at} ydre undertrykkelse dermed afløses af frivillig \dabold{selvoptimering}. Når vi konstant måler vores værd i tolv-taller, løbetider, sunde madpakker og polerede familiebilleder på sociale medier, opstår der en snigende \dabold{selvudnyttelse}. \da{Er det i grunden} muligt at gøre oprør mod en chef, når den hårdeste pisk ligger inde i ens egen bevidsthed? \da{Når man skræller overfladen væk, viser det sig, at} den uendelige jagt på det perfekte liv sjældent bringer varig ro --- tværtimod efterlader den et kronisk tomrum af \da{utilstrækkelighed}.

Denne psykologiske dynamik rammer især den unge generation. Hvor tidligere tider skabte klare grænser mellem arbejde og hvile, udviskes skellet i dag af digitale platforme og uophørlig sammenligning. Selv hvilen risikerer at blive instrumentaliseret: Man går ikke længere en tur i skoven for turens egen skyld, men for at ”lade batterierne op” til næste dags præstation. Når ethvert øjeblik skal retfærdiggøres som en investering i fremtiden, mister mennesket evnen til blot at være til stede i nuet.

Vejen ud af denne udmattelse begynder næppe med endnu en effektiv app til tidsstyring. Vi må i fællesskab genvinde modet til at sætte grænser, acceptere ufuldkommenhed og vælge langsommeligheden til. Først når værdighed ikke længere måles i resultater, kan friheden igen blive et åndehul frem for et usynligt fængsel.

\begin{transbox}
\small
\textbf{\color{NCDTeal}[Tekst A 中文精译]}\\
日常生活究竟是从何时起，变成了一个必须不断自我改良的永恒工程？闲暇时光又是从何时起，沦为了又一个用来优化身材、人脉和简历的竞技场？我们今天体验到的，归根结底究竟是真正的自由，还是只是一种更加隐蔽的内在强迫？在像丹麦这样物质基本需求大多已得到满足的福利社会里，平静按理说本该降临。然而，遭受压力、失眠与深层职业倦怠困扰的年轻人与成年人却反而日益增多。

为了理解这一谜题，韩裔德国哲学家韩炳哲指出了一个决定性的历史转向。在传统的规训社会中，人受到外部禁令、权威和明确义务的支配：“你不准”以及“你必须”。相比之下，现代绩效社会却用一种看似解放人性的声音说道：“你无所不能！”这一概念所涵盖的，并不仅仅是劳动力市场上日益严苛的要求或高中阶段对高分成绩的追逐，更是一种全面渗透的心态——在其中，个体将自己同时塑造成自己人生中的创业者与打工人。

悖论恰恰在于，外部压迫由此被自愿的“自我优化”所取代。当我们不断用满分 12 分、跑步配速、健康便当以及社交媒体上光鲜修饰的家庭照来衡量自身价值时，一种悄然蔓延的“自我剥削”便产生了。当最严厉的鞭子就握在自己脑海里时，人从根本上还能反抗老板吗？剥开表层之后便会发现，对完美生活无休止的追逐极少带来持久的安宁——恰恰相反，它留下的是一种慢性的“自我不足感”与空虚。

这种心理机制对年轻一代的冲击尤为剧烈。以往的时代曾在工作与休息之间划出清晰边界，而如今这道界线却被数字平台与无休止的横向攀比模糊了。就连休息本身也面临被工具化的危险：人们去森林散步不再是为了散步本身，而是为了给明天的绩效“充电”。当每一个瞬间都必须被证明是对未来的投资时，人就失去了纯粹安住当下的能力。

走出这种疲惫的出路，几乎不可能始于又一款高效的时间管理应用。我们必须共同重拾勇气，为生活划定边界、接纳不完美并主动选择慢下来。唯有当人的尊严不再用可量化的成绩来衡量时，自由才能重新成为喘息的空间，而非一座无形的牢笼。
\end{transbox}

\subsection*{Tekst B：\textit{Samtale i studierådgivningen: ”Man kan ikke hvile på bestilling”}（校园心理咨询对话 · 约 155 词）}

\textbf{阅读任务（理解与核对）：}\textit{Læs samtalen mellem studenterpsykolog Kasper og den universitetsstuderende Maja, og besvar spørgsmålene nedenfor.}

\begin{emailbox}{Samtale på campus｜Kasper (studenterpsykolog) \& Maja (universitetsstuderende)}
\textbf{Maja:} Kasper, jeg læste om præstationssamfundet i går, og jeg genkendte næsten alt. Selv når jeg holder fri i weekenden, får jeg dårlig samvittighed over ikke at opdatere mit CV, træne igennem eller forberede næste seminar.

\textbf{Kasper:} Det hører jeg fra studerende næsten hver eneste dag, Maja. Mange tror, at problemet blot handler om for mange arbejdstimer, men i grunden skyldes presset, at næsten alle aktiviteter bliver vurderet. Derfor har unge desperat brug for evalueringsfrie rum, hvor man kan være sammen uden at præstere.

\textbf{Maja:} Men er det ikke lidt naivt at lægge hele ansvaret over på den enkelte? Så længe uddannelsessystemet og arbejdsmarkedet belønner konstant \dabold{selvoptimering}, bliver det jo dyrt for den studerende, der vælger langsommeligheden til.

\textbf{Kasper:} Du rammer hovedet på sømmet. Individet kan ikke bryde dette mønster alene; netop derfor må både universiteter og arbejdspladser skabe en ny kultur, hvor fejl, tøven og uforpligtende pauser betragtes som naturlige.
\end{emailbox}

\textbf{Tekst B 阅读核对（Rigtigt / Forkert / Kort svar，答案见第六部分）：}
\begin{enumerate}
  \item \textit{Maja oplever, at hun uden problemer kan slappe af i weekenden uden dårlig samvittighed.} \underline{\hspace{2.2cm}}
  \item \textit{Hvad skyldes presset i grunden ifølge Kasper, og hvad har unge derfor brug for?} \underline{\hspace{2.2cm}}
  \item \textit{Hvorfor mener Maja, at den enkelte studerende ikke bare kan vælge langsommeligheden til på egen hånd?} \underline{\hspace{2.2cm}}
\end{enumerate}

\begin{tipbox}{本课发音与语调要点（Udtale \& Prosodi）}
\begin{itemize}
  \item \textbf{拉丁借词与复合抽象名词重音}：\textit{præstationssamfund} \ipa{pʰʁɛsd̥aˈɕoːˀnsˌsɑmˀfʌnˀ}（\textit{-tion-} 读 \ipa{ˈɕoːˀn}，且 \textit{samfund} 尾音 \textit{d} 不发音带 \textit{stød}）；\textit{selvoptimering} \ipa{ˈsɛlˀʌbd̥iˌmeːˀɐeŋ}；\textit{udbrændthed} \ipa{ˈuðˌbʁɛnˀdheðˀ}（软 \textit{d} 加双处 \textit{stød}）。
  \item \textbf{高频思辨词组语音细节}：\textit{i grunden} \ipa{i ˈɡʁʌnˀən}（\textit{d} 不发音！读作鼻音加 \textit{stød}）；\textit{utilstrækkelighed} \ipa{utˢilˈsd̥ʁɛɡəliˌheðˀ}（重音在 \textit{-stræk-}，结尾 \textit{g} 不发音）；\textit{paradoks} \ipa{pʰɑʁɑˈdʌɡs}（重音在最后音节 \textit{-doks}）。
  \item \textbf{连续设问的升调与节奏}：朗读以 \textit{Hvornår blev...} 开头的特殊疑问句时，丹麦语\textbf{句尾不扬反降}（核心重音落在 \textit{Hvor\textbf{når}} 和句末新信息词上）；而以 \textit{Er det i grunden...} 开头的是非问句句尾略微上扬，与前者形成极富哲学张力的语调落差。
\end{itemize}
\end{tipbox}

\begin{transbox}
\small
\textbf{\color{NCDTeal}[Tekst B 中文对照]}\\
\textbf{玛雅：}卡斯珀，我昨天读了一篇关于绩效社会的文章，里面写的一切我几乎都深有同感。哪怕在周末休息的时候，我也会因为没去更新简历、没有努力锻炼或者没提前准备下一次研讨课而感到内疚。

\textbf{卡斯珀：}玛雅，这话我几乎每天都能从学生那里听到。很多人以为问题仅仅出在学习或工作时长太长，但从根本上说，这种压力源于几乎所有活动都在被评判打分。因此，年轻人迫切需要“免于评价的空间”，在那里大家可以仅仅在一起相处而无需去表现绩效。

\textbf{玛雅：}可是把全部责任都推给个人，难道不是有点天真吗？只要教育体制和劳动力市场还在奖赏持续不断的自我优化，选择“慢下来”的学生就必然要付出高昂的代价呀。

\textbf{卡斯珀：}你一针见血地说到了关键。个体确实无法独自打破这种模式；正因如此，大学和用人单位都必须共同营造一种新文化，把犯错、犹豫以及不带功利目的的停顿视为自然之事。
\end{transbox}

% ---------------------------------------------------------------------
\ncdsubsec{三、Gloser（核心生词表）}

\begin{tipbox}{词表重点}
注意两处核心变化规律：\textit{en drivkraft}（驱动力）在变复数时元音发生变音（\textit{drivkræfter, drivkræfterne}），与 \textit{kraft $\rightarrow$ kræfter} 完全一致；强变化动词 \textit{at opstå}（产生、出现）在完成时中表示状态变化，\textbf{必须搭配助动词 \textit{er}}（\textit{er opstået}）！而 \textit{en udbrændthed}（职业倦怠）和 \textit{en utilstrækkelighed}（不足感）均为不可数抽象名词。
\end{tipbox}

\begin{glosertable}
名词 & et præstationssamfund & \textit{(præstationssamfundet, -samfund, -ene)} & 绩效社会，功绩社会 \\
名词 & et disciplinsamfund & \textit{(disciplinsamfundet, -samfund, -ene)} & 规训社会 \\
名词 & en selvoptimering & \textit{(selvoptimeringen, kun ental)} & 自我优化 \\
名词 & en selvudnyttelse & \textit{(selvudnyttelsen, kun ental)} & 自我剥削 \\
名词 & en udbrændthed & \textit{(udbrændtheden, kun ental)} & 精疲力竭，职业倦怠 \\
名词 & en utilstrækkelighed & \textit{(utilstrækkeligheden, kun ental)} & 不足感，力不从心感 \\
名词 & en tvang & \textit{(tvangen, kun ental)} & 强迫，强制力 \\
名词 & en drivkraft & \textit{(drivkraften, drivkræfter, drivkræfterne)} & 驱动力，内在动力 \\
名词 & en langsommelighed & \textit{(langsommeligheden, kun ental)} & 缓慢，慢节奏 \\
名词 & et paradoks & \textit{(paradokset, paradokser, paradokserne)} & 悖论，矛盾现象 \\
名词 & et åndehul & \textit{(åndehullet, åndehuller, åndehullerne)} & 喘息空间，避风港 \\
\midrule
动词 [har] & at præstere & \textit{(præsterer, præsterede, har præsteret)} & 表现，完成绩效 \\
动词 [har] & at optimere & \textit{(optimerer, optimerede, har optimeret)} & 优化，追求极致效率 \\
动词 [har] & at udnytte & \textit{(udnytter, udnyttede, har udnyttet)} & 剥削，利用 \\
动词 [har] & at retfærdiggøre & \textit{(retfærdiggør, -gjorde, har retfærdiggjort)} & 为……辩护，使合理化 \\
动词 [er] & at opstå & \textit{(opstår, opstod, er opstået)} & 产生，涌现，形成 \\
\midrule
形容词 & subtil & \textit{(subtilt, subtile)} & 微妙的，隐蔽难察的 \\
形容词 & kronisk & \textit{(kronisk, kroniske)} & 慢性的，长期持续的 \\
\end{glosertable}

% ---------------------------------------------------------------------
\ncdsubsec{四、Grammatik \& Stilistik（连续设问推进、概念拆解与 C1 哲学修辞）}

\begin{grammatikbox}{语法点 1：连续设问推动深度论述（\textit{Retoriske spørgsmål som argumentationsmotor}）}
在哲学与社会心理学评论中，平铺直叙容易流于干瘪；通过\textbf{排比式连续设问}（2 至 3 句紧扣的追问）直接拷问习以为常的生活前提，能瞬间确立思想深度：

\vspace{0.3em}
\renewcommand{\arraystretch}{1.24}
{\small
\begin{tabularx}{\linewidth}{@{} >{\bfseries\raggedright\arraybackslash}p{3.4cm} >{\raggedright\arraybackslash}p{4.3cm} >{\raggedright\arraybackslash}X @{}}
\toprule
\textbf{设问修辞层} & \textbf{句式功能与语序结构} & \textbf{Tekst A 高分示范} \\
\midrule
\textbf{1. 时间溯源设问}\newline \da{Hvornår blev...?} &
连用两个 \textit{Hvornår blev X til Y?}，揭示“习以为常”其实只是晚近历史产物。 &
\textit{\da{Hvornår blev} hverdagen et evigt projekt? \da{Hvornår blev} fritiden reduceret til...?} \\
\midrule
\textbf{2. 本质追问设问}\newline \da{Er det i grunden...?} &
直接质疑“自由/进步”等漂亮口号是否名实相符，准备引出概念翻转。 &
\textit{\da{Er det i grunden} ægte frihed, vi oplever i dag, eller blot en mere subtil form for tvang?} \\
\midrule
\textbf{3. 反诘悖论设问}\newline \textit{Er det... muligt, når...?} &
用自身矛盾的处境反问，让读者自己推导出“系统内无处可逃”的结论。 &
\textit{\da{Er det i grunden} muligt at gøre oprør mod en chef, når den hårdeste pisk ligger inde i én?} \\
\bottomrule
\end{tabularx}
}
\end{grammatikbox}

\begin{grammatikbox}{语法点 2：概念拆解与悖论揭示四大核心句型}
\begin{enumerate}
  \item \textbf{否定表象 + 递进深挖：\da{Begrebet dækker ikke blot over A, men over B}}：\newline
    用 \textit{dække over}（涵盖、指代）做学术界定；注意对称连接词后若接名词短语，介词 \dabold{over} \textbf{必须重复出现}：\newline
    \ding{51}\ \textit{Begrebet dækker ikke blot \dabold{over} lange arbejdsdage, men \dabold{over} en indre drivkraft.}
  \item \textbf{揭示反常机制：\da{Paradokset består i, at...}}（悖论在于……）：\newline
    动词 \textit{bestå} 必须搭配介词 \dabold{i} 再接 \textit{at}-从句（内部不倒装）：\newline
    \ding{51}\ \textit{Paradokset består \dabold{i, at} ydre undertrykkelse afløses af frivillig selvoptimering.}
  \item \textbf{剥离表象隐喻：\da{Når man skræller overfladen væk, viser det sig, at...}}（剥开表层就会发现……）：\newline
    句首长条件时间从句占据第 1 位，因此主句必须紧跟变位动词 \dabold{viser}（V2 倒装）：\newline
    \ding{51}\ \textit{Når man skræller overfladen væk, \dabold{viser det sig}, at jagten skaber utilstrækkelighed.}
  \item \textbf{历史对照界定：\da{I det klassiske A ... Det moderne B taler derimod...}}：\newline
    通过 \textit{derimod}（相比之下，置于主句动词后）把传统规训与现代绩效并置对比。
\end{enumerate}
\end{grammatikbox}

\begin{grammatikbox}{语法点 3：哲学评论常见中丹语序与搭配误区}
\begin{enumerate}
  \item \textbf{\ding{172} \textit{i grunden} 放主句与放从句的位置差异}：在主句是非问句里放在主语后（\textit{Er det \da{i grunden} frihed?}）；在普通主句中置于动词后（\textit{Det handler \da{i grunden} om vurdering.}）；在 \textit{at}-从句中置于动词前（\textit{...at friheden \da{i grunden} er en ny tvang.}）。
  \item \textbf{\ding{173} \textit{bestå i} vs \textit{bestå af}}：\textit{bestå af} 指“由具体零部件/成员组成”；阐述抽象道理、核心原因或悖论所在时，\textbf{只能用 \da{bestå i}}（不能误作 \ding{55}\ \textit{Paradokset består af, at...}）。
  \item \textbf{\ding{174} \textit{at opstå} 与 \textit{at forsvinde} 必配助动词 \textit{er}}：凡表示状态“产生/消失”的不可及物变化动词，完成时均用 \textit{er}，不可误作 \ding{55}\ \textit{der har opstået en ny træthed} $\rightarrow$ \ding{51}\ \textit{der \da{er opstået} en ny træthed}。
  \item \textbf{\ding{175} 分清 \textit{i stedet for} 与 \textit{frem for}}：C1 抽象价值抉择中，表达“宁可要 A 而不要 B（……而非……）”常用 \da{frem for}：\textit{et åndehul \da{frem for} et usynligt fængsel}。
\end{enumerate}
\end{grammatikbox}

% ---------------------------------------------------------------------
\ncdsubsec{五、Mønsterøvelser（句型与修辞操练）}

\textbf{练习 A：连续设问、概念拆解与绩效社会核心词精准填空（每空填入最准确的 1 个丹麦语词）}
\begin{enumerate}
  \item \underline{\hspace{2.2cm}} blev hverdagen et projekt, der hele tiden skal forbedres og optimeres?
  \item Er det i \underline{\hspace{2.2cm}} ægte frihed, når man aldrig tør holde fri uden dårlig samvittighed?
  \item Begrebet dækker ikke blot over lange arbejdstimer, men \underline{\hspace{2.2cm}} en dyb indre drivkraft.
  \item Paradokset består \underline{\hspace{2.2cm}}, at borgerne frivilligt overvåger og presser sig selv.
  \item Når man skræller overfladen væk, \underline{\hspace{2.2cm}} det sig, at jagten på perfektion skaber tomhed.
  \item I det moderne præstationssamfund er der \underline{\hspace{2.2cm}} en udbredt følelse af utilstrækkelighed. (opstå，完成分词)
  \item Konstant \underline{\hspace{2.2cm}} på sociale medier fører let til stress og udbrændthed hos unge. (自我优化)
  \item Unge har brug for et mentalt \underline{\hspace{2.2cm}}, hvor de ikke hele tiden bliver målt og vurderet. (喘息空间)
\end{enumerate}

\textbf{练习 B：哲学思辨句法纠错与修辞升级（修正搭配/语序错误，或按提示升级）}
\begin{enumerate}
  \item 【悖论介词纠错】\textit{Paradokset består af, at vi har mere fritid, men føler os mere stressede end før.} $\rightarrow$
  \item 【助动词与从句语序纠错】\textit{Artiklen forklarer, at der har opstået en ny form for selvudnyttelse blandt studerende.} $\rightarrow$
  \item 【对称介词与平行纠错】\textit{Udbrændthed dækker ikke blot over træthed i kroppen, men også en følelse af manglende værdi.} $\rightarrow$
  \item 【前置时间从句 V2 纠错】\textit{Når man skræller de polerede billeder væk, det viser sig, at mange kæmper med angst.} $\rightarrow$
  \item 【连续设问修辞升级】将平淡陈述改写为“\textit{Hvornår blev...}”修辞设问：\textit{Vores fritid er desværre blevet til en slags arbejde, hvor man skal optimere sin sundhed.} $\rightarrow$
  \item 【概念对照升级】用 \textit{I modsætning til ... er ... præget af ...} 整合句意：\textit{Disciplinsamfundet brugte ydre forbud. Præstationssamfundet bruger indre selvoptimering.} $\rightarrow$
\end{enumerate}

\textbf{练习 C：中译丹实战（运用连续设问与概念拆解表达，精准锤炼 C1 语感）}
\begin{enumerate}
  \item 日常生活是从何时起变成一项必须不断优化的工程的？
  \item 我们今天体验到的，归根结底究竟是自由还是一种隐蔽的强迫？
  \item 这一概念不仅涵盖高分要求，也涵盖持续不断的自我剥削。
  \item 悖论在于，外部指令已被现代人的自愿自我优化所取代。
  \item 剥开表层就会发现，对完美生活的追逐反而滋生出慢性不足感。
  \item 当任何休息都必须被证明是合理的，人们便失去了当下的喘息空间。
\end{enumerate}

% ---------------------------------------------------------------------
\ncdsubsec{六、PD3 / Studieprøven 实战微任务与参考答案}

\begin{emailbox}{PD3 / Studieprøven 写作微任务（Reflekterende essay · 200--250+ ord）}
\textbf{Situation}: En ny undersøgelse viser, at næsten hver tredje unge dansker oplever dagligt præstationspres i forbindelse med studier, fritid og sociale medier. På en ungdomsuddannelse diskuterer man nu, om man skal afskaffe karakterer i det første år for at dæmpe presset.\\[3pt]
\textbf{Opgave}: Skriv et reflekterende debatindlæg på \textbf{200--250 ord}, hvor du analyserer årsagerne til præstationspresset og tager stilling til forslaget. Brug mindst \textbf{to retoriske spørgsmål} til at drive argumentationen frem samt \textbf{tre} begrebsanalytiske udtryk fra lektionen (\textit{præstationssamfund, selvoptimering, paradokset består i, at..., i grunden, utilstrækkelighed, udbrændthed}).
\end{emailbox}

\begin{answerbox}
\textbf{Før du læser 参考回答}
\begin{enumerate}
  \item \textit{Ifølge Byung-Chul Han blev mennesket i disciplinsamfundet styret af ydre forbud og pligter, mens præstationssamfundet lokker med ”Du kan alt”, så individet selv bliver både chef og arbejder. Når værdien hele tiden måles gennem selvoptimering, opstår en kronisk følelse af utilstrækkelighed og udbrændthed.}
  \item \textit{Kasper mener, at unge har brug for evalueringsfrie rum, hvor de blot kan være sammen uden at blive vurderet. Maja påpeger imidlertid, at den enkelte studerende betaler en høj pris for at sige nej, så længe uddannelsessystemet og arbejdsmarkedet fortsat belønner konstant selvoptimering.}
\end{enumerate}

\textbf{Tekst B 阅读核对答案}
\begin{enumerate}
  \item \da{Forkert}（玛雅一到周末反而因为没有完善简历或预习课程而感到内疚：\textit{får jeg dårlig samvittighed}）
  \item \da{At næsten alle aktiviteter bliver vurderet; evalueringsfrie rum.}（根本原因在于几乎一切活动都在被评价打分，因此需要免于评价的空间）
  \item \da{Fordi uddannelsessystemet og arbejdsmarkedet belønner konstant selvoptimering.}（因为教育体制和劳动力市场仍在奖赏持续的自我优化）
\end{enumerate}

\textbf{练习 A 答案}
\begin{enumerate}
  \item \da{Hvornår}（时间溯源连续设问核心词：\textit{Hvornår blev...}）
  \item \da{grunden}（本质追问副词短语：\textit{i grunden} 归根结底、从根本上说）
  \item \da{over}（对称复现介词：\textit{dækker ikke blot over A, men over B}）
  \item \da{i}（固定搭配：\textit{bestå i, at...} 在于……，切勿误用 \textit{af}）
  \item \da{viser}（前置时间从句后触发主句 V2 倒装：\textit{viser det sig}）
  \item \da{opstået}（状态变化动词 \textit{at opstå} 与助动词 \textit{er} 构成现在完成时）
  \item \da{selvoptimering}（课内核心哲学心理概念名词：自我优化）
  \item \da{åndehul}（名词：喘息空间，避风港）
\end{enumerate}

\textbf{练习 B 答案（参考）}
\begin{enumerate}
  \item \textit{Paradokset består \da{i}, at vi har mere fritid, men føler os mere stressede end før.}（阐述“悖论在于……”须用 \textit{bestå i}，不可误用 \textit{bestå af}）
  \item \textit{Artiklen forklarer, at der \da{er opstået} en ny form for selvudnyttelse blandt studerende.}（\textit{opstå} 属状态变化动词，完成时必须搭配助动词 \textit{er}）
  \item \textit{Udbrændthed dækker ikke blot \da{over} træthed i kroppen, men også \da{over} en følelse af manglende værdi.}（注意 \textit{ikke blot ... men også ...} 并列两个名词短语时须对称复现介词 \textit{over}）
  \item \textit{Når man skræller de polerede billeder væk, \da{viser det sig}, at mange kæmper med angst.}（前置从句占第 1 位，主句变位动词 \textit{viser} 必须紧跟实现 V2 倒装）
  \item \textit{\da{Hvornår blev} vores fritid reduceret til et arbejde, hvor man skal optimere sin sundhed?}（借助连续设问将平淡叙述升级为尖锐哲学追问）
  \item \textit{\da{I modsætning til} disciplinsamfundet, der brugte ydre forbud, er præstationssamfundet \da{præget af} indre selvoptimering.}（运用 \textit{I modsætning til} 与关系从句完成概念对照拆解）
\end{enumerate}

\textbf{练习 C 答案}
\begin{enumerate}
  \item \textit{Hvornår blev hverdagen til et projekt, som man hele tiden skal optimere?}
  \item \textit{Er det i grunden frihed eller en subtil form for tvang, vi oplever?}
  \item \textit{Begrebet dækker ikke blot over høje karakterkrav, men også over konstant selvudnyttelse.}
  \item \textit{Paradokset består i, at ydre pligter er blevet erstattet af frivillig selvoptimering.}
  \item \textit{Når man skræller overfladen væk, nærer jagten på perfektion en kronisk utilstrækkelighed.}
  \item \textit{Hvis enhver pause skal retfærdiggøres, mister vi vores nødvendige åndehul i nuet.}
\end{enumerate}

\textbf{PD3 / Studieprøven 高分范文（约 235 词，不含标题）}\\
\textit{\textbf{Karakterfrihed er kun første skridt ud af præstationsfælden}}

\textit{\da{Hvornår blev} ungdomsårene forvandlet fra en tid med nysgerrighed, vildveje og fri eksperimenteren til et stramt optimeringsprojekt? \da{Er det i grunden} mangel på viljestyrke, der plager dagens unge, eller reagerer de blot fornuftigt på et grænseløst ydre og indre pres? Når næsten hver tredje unge dansker oplever dagligt stress i \dabold{præstationssamfundet}, peger forslaget om et karakterfrit første studieår i den rigtige retning --- men det løser næppe hele problemet alene.}

\textit{\da{Begrebet præstationspres dækker nemlig ikke blot over} frygten for en lav karakter på eksamensbeviset, \da{men over} en gennemgribende kultur, hvor selv fritid, venskaber og krop underkastes konstant \dabold{selvoptimering}. \da{Paradokset består i, at} jo mere frihed de unge formelt har fået til at forme deres egen fremtid, desto hårdere dømmer og overvåger de sig selv gennem sociale medier og urealistiske idealer. \da{Når man skræller overfladen væk, viser det sig, at} den udbredte \dabold{udbrændthed} og træthed ofte udspringer af en kronisk følelse af \da{utilstrækkelighed}.}

\textit{Et karakterfrit første år kan skabe et værdifuldt og nødvendigt åndehul, hvor fejl ikke straks straffes, og hvor samtale, leg og faglig fordybelse igen får plads uden målestok. Hvis vi imidlertid fortsat måler menneskelig værdi på hastighed, effektivitet og fejlfrie resultater i resten af uddannelsessystemet og arbejdslivet, flytter presset blot videre til andet år. Ægte forandring kræver derfor både konkrete reformer i skolen, evalueringsfrie rum i hverdagen og et modigt opgør med selve præstationskulturen uden for klasseværelset.}

\small\textbf{名师点评}：正文（不含标题）约 235 词，堪称哲学心理学思辨佳作。开篇两连设问（\textit{Hvornår blev...} 与 \textit{Er det i grunden...}）直击问题本质；中段以 \textit{Begrebet ... dækker ikke blot over..., men over...} 拆解“绩效压力”的深层结构，再用 \textit{Paradokset består i, at...} 与 \textit{Når man skræller overfladen væk, viser det sig, at...} 揭示年轻一代越自由越疲惫的深层悖论；末段辩证评估“高一取消打分制”的限度，既有理论高度又有教育现实感！
\end{answerbox}